% !TeX root = ../main.tex
\section{Discussion}
\label{sec:discussion}

\noindent\textbf{Partial Task 4.1 draft.} The computational answers below
use the approved Results; the expert-assessment answer remains unresolved.

\subsection{Answers to the Research Questions}

\emph{RQ1---Observed applicability classifications.}
Within the sampled requirements and catalog, the LLM predominantly
returned insufficient information: \StudyInsufficientInformationCount{}
of \StudyAssessmentCount{} pairs
(\RqOneInsufficientInformationPercent{}\%).
It identified \StudyApplicableCount{} pairs as applicable
(\RqOneApplicablePercent{}\%), spanning \StudyPositiveStoryCount{}
of \StudyStoryCount{} stories; \StudyNotApplicableCount{} pairs were
classified as not applicable (Table~\ref{tab:applicability-results}).
Thus, the output is a bounded set of story-specific pattern candidates,
not a universal recommendation list. The dominant unresolved category
limits what can be concluded about catalog applicability from this run:
it does not establish that those patterns are unsuitable or that the
requirements are defective. Likewise, positive classifications are model
judgments, not validated architectural choices. Differences across stories
describe this sample; explaining them would require evidence beyond the
reported label frequencies.

\emph{RQ2---Overlap with the frozen retrieval baseline.}
The methods' outputs intersect on \StudyRagOverlapCount{} of the
\StudyApplicableCount{} LLM-positive pairs
(\RqTwoOverlapPercent{}\%), across \StudyRagOverlapStoryCount{} stories;
\StudyRagOutsideCount{} positive pairs lie outside the same-story final
Top-5 (Table~\ref{tab:rag-comparison}).
This answers a candidate-membership question: the final Top-5 does not
contain most pairs labeled applicable by the exhaustive assessment.
It does not establish which method makes better architectural selections.
The selected-pair set is not an independent reference, and the two
protocols differ in input language and catalog representation.
Consequently, non-overlap cannot be interpreted as retrieval error,
additional validated discoveries, or superiority of exhaustive assessment.

\emph{RQ3---Expert assessment pending.}
\ExpertResultsUnresolved{}. Table~\ref{tab:expert-evaluation-results}
is an unresolved reporting shell, not evidence of expert judgments.
The expert judgments remain unknown. The complete, traceable ADR outputs
(Section~\ref{sec:adr-results}) establish generation completion, not
professional acceptance or semantic quality.
Whether experts corroborate, qualify or challenge the recommendations
remains undetermined. Revisit this answer and the combined interpretation
after Phase 3 verifies ratings, coverage and eligible comments;
no outcome is assumed.

\subsection{Interpreting Exhaustive Assessment and Retrieval}
\label{sec:llm-rag-discussion}

The two protocols have different decision boundaries.
Exhaustive assessment evaluates each story--pattern pair against the
supplied requirements and returns an applicability category with
supporting evidence for positive classifications
(Section~\ref{sec:applicability-experiment}).
The retrieval baseline instead ranks catalog documents by their
relationship to a story query and retains a fixed-size candidate set
(Section~\ref{sec:rag-baseline}). A ranking position is not an
applicability judgment: even a story with no LLM-positive pair receives
retrieval candidates. Conversely, exhaustive assessment is not constrained
to return a fixed number of positives. Their intersection therefore
describes agreement in candidate membership, not agreement on a common
classification task.

Table~\ref{tab:rag-comparison} shows that the frozen final Top-5 omits
most LLM-positive pairs. This observation neither proves that retrieval
misses appropriate patterns nor validates the additional LLM selections.
English retrieval queries, Portuguese assessment inputs and unequal
catalog representations prevent attributing the difference solely to
the selection mechanism. The current comparison also does not establish
why individual candidates rank outside the cutoff. Architectural quality
and relevance require evidence independent of either output set; expert
assessment of selected pairs alone cannot establish recall across the
unassessed alternatives.

\emph{Future design hypothesis---not evaluated.}
Retrieval could serve as a candidate-generation stage before pairwise
LLM assessment, potentially reducing the number of pairs sent for
assessment. Such filtering would also exclude the current LLM-positive
pairs outside the cutoff from consideration; the architectural importance
of those exclusions is unknown. No hybrid pipeline was executed, and
end-to-end cost savings, latency and quality effects are
\textbf{not evaluated}, not demonstrated benefits.
A controlled follow-up should align input language and catalog fields,
vary the candidate budget, account for retrieval/reranking overhead
alongside generation costs, and independently assess candidates both
inside and outside the shortlist. Until those trade-offs are measured,
retrieval-first assessment remains a proposed alternative rather than
a validated replacement for exhaustive assessment.
